% bab1.tex — dari BAB_I_PENDAHULUAN_POLISHED.md mulai ## A.

\phantomsection
\section*{BAB I --- PENDAHULUAN}
\addcontentsline{toc}{section}{BAB I --- PENDAHULUAN}

\phantomsection
\subsection*{A. Latar Belakang}
\addcontentsline{toc}{subsection}{A. Latar Belakang}

Hadis, selaku sumber kedua ajaran Islam sesudah Al-Qur'an, menduduki posisi pivotal dalam ranah hukum, teologi, maupun etika Islam lewat perannya sebagai \textit{bayan} (klarifikasi), \textit{takhsis} (pembatasan makna umum), serta legislasi otonom atas kasus yang tidak dirinci Al-Qur'an.\footnote{A. Musadad, "Articulating the prophetic authority: Some notes on the nature of hadith collection in al-Muwatta' and musannaf literature," \\textit{Jurnal Studi Ilmu-Ilmu Al-Qur'an dan Hadis}, vol. 27, no. 1, hlm. 351--373, 2026; N. M. A. M. Razak et al., "Evolution of hadith textual criticism: From classical to contemporary approaches," \\textit{Global Journal Al-Thaqafah}, hlm. 15--29, 2025.} Wibawa kenabian itu dikristalkan lewat artikulasi panjang himpunan perdana semisal \textit{al-Muwatta'} beserta \textit{Musannaf}, yang memperlihatkan bagaimana otoritas Nabi dibangun secara metodis oleh ulama generasi pertama.\footnote{Musadad, "Articulating the Prophetic Authority."; T. N. S. Tanjung and Z. Fadhlurrahman, "The role of tadwin hadith in the development of Islamic historiography," \\textit{Istifham: Journal of Islamic Studies}, vol. 3, no. 3, hlm. 188--201, 2025.} Posisi itu kian menguat dalam fikih masa kini, sebab kritik hadis yang tajam terbukti berkonsekuensi hukum langsung --- umpamanya pada wacana kewarisan Islam, di mana pengujian \textit{matan} dan \textit{sanad} memutus keabsahan dalil.\footnote{M. R. Rozikin, "Al-fara'id is half of knowledge: A critical hadith analysis and its juridical implications in Islamic inheritance discourse," \\textit{Cogent Arts and Humanities}, vol. 13, no. 1, 2026; K. H. Bin Jamil, "Between traditionalising and futuring: Applying the broader maqasid paradigm to hadith studies," \\textit{Asia Pacific Journal of Educators and Education}, vol. 39, no. 2, hlm. 183--196, 2024.} Pada level epistemologi, hadis menjelma titik temu antara khazanah klasik dan kerangka \textit{maqasid}, sampai-sampai telaahnya melebar dari sekadar verifikasi tekstual menuju orientasi kemaslahatan syariat.\footnote{Bin Jamil, "Between Traditionalising and Futuring."; Razak, "Evolution of Hadith Textual Criticism.".} Konsekuensinya, penguasaan hadis secara utuh meniscayakan penelusuran serentak atas dimensi kesejarahan, metodologi, dan kontekstualitasnya, karena setiap klaim normatif yang dinisbatkan kepada Nabi berdampak langsung pada amalan hukum serta moral umat.\footnote{Rozikin, "Al-fara'id Is Half of Knowledge."; Musadad, "Articulating the Prophetic Authority.".}

Dari sudut kesejarahan, lintasan hadis terhampar sejak era kenabian hingga zaman modern lewat metamorfosis bertingkat: berawal dari kelisanan, catatan personal (\textit{sahifah}), pembukuan tematik, hingga penetapan kitab induk.\footnote{Tanjung, "The Role of Tadwin Hadith."; M. Lutfianto et al., "History of hadith writing in the precodification period in the prospective of 'Ajjaj al-Khatib," \\textit{Al-Thiqah: Jurnal Ilmu Keislaman}, vol. 7, no. 2, hlm. 308, 2024.} Di masa Nabi, periwayatan berjalan lewat pelbagai wahana --- hafalan, keteladanan praktis, pembelajaran majelis, dan dokumentasi dini selaku penguat ingatan --- sehingga periode tersebut tidak tepat disederhanakan menjadi sepenuhnya oral ataupun sepenuhnya tertulis.\footnote{H. Rashwan, "Hadith as oral literature through early Islamic literary criticism," \\textit{Studia Islamica}, vol. 119, no. 1, hlm. 34--69, 2024; S. Watukila, "Historiografi hadis di masa mutaqaddimin dalam pandangan Mustafa Azami," \\textit{Amsal Al-Qur'an: Jurnal Al-Qur'an dan Hadis}, vol. 1, no. 3, hlm. 221--235, 2024.} Selepas wafat beliau, hijrah para sahabat ke Syam, Irak, serta Mesir memunculkan transmisi yang bercorak kedaerahan, sementara mobilitas intelektual (\textit{rihlah}) mengarahkan para sahabat menjadi rujukan otoritatif di tiap kawasan.\footnote{R. E. Gil, "Tracing the rihla in buldan books: Ibn al-Faqih's muhaddith resources," \\textit{Hitit Theology Journal}, vol. 23, hlm. 140--154, 2024; H. B. Jaiyeoba and N. M. Osmani, "Hadith preservation: Techniques and contemporary efforts," \\textit{Journal of Fatwa Management and Research}, vol. 29, no. 3, hlm. 31--45, 2024.} Memasuki generasi tabi'in, jejaring guru-murid antar-kawasan kian solid, simpul-simpul keilmuan regional pun terbentuk, dan \textit{isnad} dimantapkan sebagai kaidah pembuktian yang kian berkesadaran metodis.\footnote{Watukila, "Historiografi Hadis Di Masa Mutaqaddimin."; Lutfianto, "History of Hadith Writing.".} Puncaknya pada era ulama klasik, kompilasi semisal \textit{al-Muwatta'}, \textit{al-Musannaf}, dan \textit{Sahihayn} memindahkan tumpuan otoritas dari figur perawi menuju kitab kanon beserta perangkat kritiknya yang telah mapan.\footnote{Musadad, "Articulating the Prophetic Authority."; S. M. S. Obeid and F. S. Hussein, "Al-Bukhari's sources," \\textit{Jordan Journal for History and Archaeology}, vol. 17, no. 3, hlm. 44--62, 2023.} Telaah atas penulisan hadis pra-kodifikasi melalui lensa 'Ajjaj al-Khatib menegaskan bahwa peralihan hafalan menuju penulisan berlangsung secara gradual, dengan \textit{qira'ah}, \textit{sama'}, serta audisi berperan selaku jembatan menuju pembukuan menyeluruh.\footnote{Lutfianto, "History of Hadith Writing."; Tanjung, "The Role of Tadwin Hadith.".}

Pada rentang lintasan itu mencuat problem akut berupa fabrikasi hadis (\textit{maudu'}) --- kabar buatan yang diklaim bersumber dari Nabi, padahal beliau sama sekali tidak pernah menyabdakannya.\footnote{M. A. M. Muhamed Ali et al., "Al-jarh wa al-ta'dil (criticism and praise): Its significance in the science of hadith," \\textit{Mediterranean Journal of Social Sciences}, vol. 6, no. 2S1, hlm. 284, 2015; Wasman et al., "A critical approach to prophetic traditions: Contextual criticism in understanding hadith," \\textit{Al-Jami'ah: Journal of Islamic Studies}, vol. 61, no. 1, hlm. 1--17, 2023.} Batas struktural antara zaman kemurnian periwayatan dan zaman pemalsuan berpangkal pada gejolak politik seusai terbunuhnya Khalifah Usman dan meletusnya \textit{al-fitnah al-kubra}, yang menelurkan kubu politik-teologis yang saling berebut legitimasi.\footnote{H. H. Liew, "The caliphate will last for thirty years: Polemic and political thought in the afterlife of a prophetic hadith," \\textit{Journal of Islamic Studies}, vol. 36, no. 1, hlm. 38--82, 2025; I.-W. Su, "'Ali b. al-Madini (161-234/778-849): A critical reconstruction of his biography and evaluation of his contribution to hadith criticism," \\textit{Journal of Islamic Studies}, vol. 33, no. 1, hlm. 1--34, 2021.} Riset mengenai \textit{afterlife} hadis kenabian seputar khilafah tiga puluh tahun memperlihatkan bahwa satu \textit{matan} yang sama dapat dimaknai-ulang guna menyangga narasi kekuasaan yang bertolak belakang.\footnote{Liew, "The Caliphate Will Last for Thirty Years."; S. Kara, "Debating the origins: The sanctity of Madina in hadith narratives," \\textit{Journal of Islamic Studies}, vol. 37, no. 2, hlm. 163--207, 2026.} Di luar arena politik, pecahnya teologi (Murji'ah, Qadariyah, Jabariyah), rivalitas mazhab fikih, serta ambisi popularitas dan imbalan materiil para tukang cerita (\textit{qussas}) turut menyuburkan industri pemalsuan pada abad kedua--ketiga Hijriah.\footnote{Wasman, "A Critical Approach to Prophetic Traditions."; U. M. Noor, "Muqaddimah of Ibn al-Salah and the revival of hadith studies in the Mamluk era," \\textit{Al-Bayan: Journal of Qur'an and Hadith Studies}, vol. 21, no. 2, hlm. 271--297, 2023.} Pemetaan jaringan perawi \textit{Sahih Bukhari} menampakkan betapa rumitnya bangunan transmisi itu, sampai-sampai penyimpangan relasional di dalam jejaring \textit{sanad} dapat dibaca sebagai petunjuk adanya distorsi atau sisipan.\footnote{T. Alam and J. Schneider, "Social network analysis of hadith narrators from Sahih Bukhari," \\textit{Proceedings of 2020 7th IEEE International Conference on Behavioural and Social Computing (BESC)}, hlm. 1--6, 2020; M. A. Mosa, "Synergizing structure and semantics: A knowledge graph-transformer framework for narrator disambiguation in hadith networks," \\textit{Digital Scholarship in the Humanities}, 2025.} Dengan kata lain, fabrikasi hadis bukan sekadar penyimpangan literatur pinggiran, melainkan krisis teologis berskala besar yang merongrong otentisitas ajaran dari dalam dan mendesak jawaban metodologis yang sistematis.\footnote{Muhamed Ali, "Al-jarh Wa Al-ta'dil (criticism and Praise)."; Wasman, "A Critical Approach to Prophetic Traditions.".} Jawaban ulama pun berwujud bangunan penyelamatan yang berpijak pada kerangka \textit{al-jarh wa al-ta'dil}, \textit{takhrij}, serta kritik \textit{matan}.\footnote{Muhamed Ali, "Al-jarh Wa Al-ta'dil (criticism and Praise)."; S. F. Abbas and M. A. Rawabdeh, "The concept of munkar in al-Dhahabi's critique of al-Hakim's hadith authentication in al-Mustadrak," \\textit{Jurnal Studi Ilmu-Ilmu Al-Qur'an dan Hadis}, vol. 26, no. 2, hlm. 545--568, 2025.} Arti penting \textit{al-jarh wa al-ta'dil} bersumber dari fungsinya selaku penyaring integritas (\textit{'adalah}) serta kekuatan hafalan (\textit{dabt}) perawi, sehingga hanya periwayatan perawi \textit{thiqah} yang diterima sebagai argumen.\footnote{Muhamed Ali, "Al-jarh Wa Al-ta'dil (criticism and Praise)."; Su, "'ali B. Al-madini (161-234/778-849).".} Rekonstruksi atas sumbangsih 'Ali b. al-Madini (161--234 H/778--849 M) memperlihatkan bahwa kritik perawi pada generasi dini telah menanam fondasi pembakuan ilmu \textit{'ilal} serta \textit{rijal} yang kelak diteruskan Ibn al-Salah beserta generasi Mamluk.\footnote{Su, "'ali B. Al-madini (161-234/778-849)."; Noor, "Muqaddimah of Ibn Al-salah.".} Pada ranah \textit{matan}, kupasan kritis Ibn Kasir terhadap sebagian hadis \textit{Sahihayn} bertema \textit{sirah} di dalam \textit{al-Bidaya wa al-Nihaya} memperlihatkan bahwa otoritas kanonik pun tetap terbuka bagi telaah isi berbasis keselarasan historis-teologis.\footnote{M. A. Calgan, "Analysis of Ibn Kathir's content criticism of some of the Sahihayn hadiths in the sirah section of al-Bidaya wa al-Nihaya," \\textit{Cumhuriyet Ilahiyat Dergisi}, vol. 28, no. 1, hlm. 303--324, 2024; N. Ardig et al., "Genre analysis and religious texts: A methodological model of hadith commentary," \\textit{Journal of Islamic Studies}, vol. 36, no. 2, hlm. 163--218, 2025.} Kajian atas konsep \textit{munkar} pada telaah kritis al-Dzahabi mengenai autentikasi al-Hakim (\textit{al-Mustadrak}) mempertegas bahwa autentikasi bukan putusan tunggal, melainkan dialektika antar-muhaddis yang senantiasa diuji.\footnote{Abbas, "The Concept of Munkar."; Obeid, "Al-bukhari's Sources.".} Lewat penyusunan kitab-kitab \textit{maudu'at}, pembakuan \textit{muqaddimah} ilmu hadis, serta penyaringan sumber yang ketat menurut teladan al-Bukhari, ulama klasik menegakkan daya tahan epistemologis yang melindungi korpus dari penyusupan lanjutan.\footnote{Obeid, "Al-bukhari's Sources."; Noor, "Muqaddimah of Ibn Al-salah.".}

Memasuki abad kedua puluh satu, hadis berhadapan dengan guncangan global yang dipicu revolusi digital, jejaring sosial, dan kecerdasan artifisial generatif, sehingga transmisi beralih dari \textit{sanad} yang otoritatif menuju otoritas algoritmik dalam ranah siber.\footnote{M. A. Abdulrahman, "The future of hadith studies in the digital age: Opportunities and challenges," \\textit{Journal of Ecohumanism}, vol. 3, no. 8, 2024; A. M. B. Matin Bin Salman, "Reconstructing hadith discourse in the digital age: From text to discourse," \\textit{Journal of Ecohumanism}, vol. 4, no. 1, hlm. 1--11, 2025.} Hadis kini bukan lagi sekadar teks yang disalurkan, melainkan wacana yang dihasilkan, diedarkan, dan dimaknai-ulang dalam ekosistem digital yang terdesentralisasi.\footnote{Matin Bin Salman, "Reconstructing Hadith Discourse in the Digital Age."; Abdulrahman, "The Future of Hadith Studies in the Digital Age.".} Gejala viralitas tanpa validitas menjelma problem gawat manakala cuplikan hadis disebarluaskan secara masif tanpa \textit{takhrij} yang layak, sehingga digitalisasi \textit{takhrij} berubah menjadi keperluan mendesak demi merapatkan laju sirkulasi dengan kecermatan verifikasi.\footnote{Z. Abdullah et al., "Viralitas dan validitas: Digitalisasi takhrij hadis di era media sosial," \\textit{Dirayah: Jurnal Ilmu Hadis}, vol. 6, no. 2, hlm. 381--390, 2026; T. Supriyadi et al., "Action research in hadith literacy: A reflection of hadith learning in the digital age," \\textit{International Journal of Learning, Teaching and Educational Research}, vol. 19, no. 5, hlm. 99--124, 2020.} Riset \textit{living hadith} pada era kecerdasan artifisial menunjukkan bahwa mahasiswa masa kini menjadikan \textit{ChatGPT} sebagai rujukan belajar hadis --- kebiasaan yang merentangkan akses sekaligus menebar bahaya disinformasi keagamaan tanpa bekal literasi kritis yang cukup.\footnote{H. Syukur et al., "Living hadith in the age of artificial intelligence: A phenomenological study of the use of ChatGPT as a source for studying hadith among college students," \\textit{Asian Journal of Islamic Studies and Da'wah}, vol. 4, no. 3, hlm. 189--204, 2026; Supriyadi, "Action Research in Hadith Literacy.".} Masa depan kajian hadis diwarnai tarik-menarik antara peluang pembakuan korpus digital dan terkikisnya otoritas tradisional; implikasinya, verifikasi otomatis yang dipadukan literasi hadis digital menjadi kebutuhan prioritas.\footnote{Abdulrahman, "The Future of Hadith Studies in the Digital Age."; Bin Jamil, "Between Traditionalising and Futuring.".}

\textit{State of the art} rentang 2015--2026 memperlihatkan tiga kutub yang saling menggenapi: pemantapan kritik \textit{matan} mutakhir, komputasionalisasi korpus, dan jawaban atas guncangan digital.\footnote{Razak, "Evolution of Hadith Textual Criticism."; A. M. Azmi et al., "Computational and natural language processing based studies of hadith literature: A survey," \\textit{Artificial Intelligence Review}, vol. 52, no. 2, hlm. 1369--1414, 2019; Abdulrahman, "The Future of Hadith Studies in the Digital Age.".} Pertama, perkembangan kritik teks mencatat migrasi dari hegemoni \textit{sanad} ke arah integrasi \textit{sanad-cum-matan} plus pembacaan kontekstual yang peka \textit{asbab al-wurud}.\footnote{Razak, "Evolution of Hadith Textual Criticism."; Wasman, "A Critical Approach to Prophetic Traditions.".} Telaah perbandingan dampak modernitas terhadap kritik \textit{matan} pada \textit{Fath al-Bari} karya Ibn Hajar serta \textit{Majallat al-Manar} karya Rasyid Rida menampakkan dua kutub, tradisionis dan reformis, dalam menanggapi hadis bertema gejala alam.\footnote{A. A. R. Al-Imam, "The influence of modernity on matn criticism: A comparative study of traditionist and reformist approaches in Ibn Hajar's Fath al-Bari and Rashid Rida's Majallat al-Manar on hadiths of natural phenomena," \\textit{Journal of Islamic Thought and Civilization}, vol. 16, no. 1, hlm. 01--17, 2026; Razak, "Evolution of Hadith Textual Criticism.".} Kedua, komputasionalisasi melahirkan model \textit{pretrained} bagi pengenal hadis palsu, pemodelan corak kenabian lewat \textit{AraBERT}, dan kerangka graf pengetahuan-\textit{transformer} untuk penguraian ambiguitas narator.\footnote{J. Alghamdi et al., "Pretrained models against traditional machine learning for detecting fake hadith," \\textit{Electronics}, vol. 14, no. 17, hlm. 3484, 2025; M. Shaaban et al., "Prophetic authorial style modeling for detecting fabricated hadiths using AraBERT," \\textit{Journal of Computing & Biomedical Informatics}, vol. 10, no. 2, 2026; Mosa, "Synergizing Structure and Semantics.".} Survei riset berbasis pemrosesan bahasa alami beserta penilaian fitur kebahasaan melalui \textit{machine learning} mengukuhkan capaian akurasi prediksi yang berarti pada pendekatan komputasional, walaupun validasi ulama tetap tak tergantikan.\footnote{Azmi, "Computational and Natural Language Processing Based Studies of Hadith Literature."; E. Mohamed and R. Sarwar, "Linguistic features evaluation for hadith authenticity through automatic machine learning," \\textit{Digital Scholarship in the Humanities}, vol. 37, no. 3, hlm. 830--843, 2022.} Ketiga, ramalan otentisitas lewat telaah sentimen, komparasi \textit{deep learning} lawan kompresi lawan pengklasifikasi klasik, dan ontologi \textit{SemanticHadith} memperlihatkan kematangan prasarana semantik verifikasi berskala luas.\footnote{F. Haque et al., "Hadith authenticity prediction using sentiment analysis and machine learning," \\textit{2020 IEEE 14th International Conference on Application of Information and Communication Technologies (AICT)}, hlm. 1--6, 2020; T. Tarmom et al., "Deep learning vs compression-based vs traditional machine learning classifiers to detect hadith authenticity," \\textit{Communications in Computer and Information Science}, hlm. 206--222, 2022; A. B. Kamran et al., "SemanticHadith: An ontology-driven knowledge graph for the hadith corpus," \\textit{Journal of Web Semantics}, vol. 78, hlm. 100797, 2023.} Kedudukan riset ini ialah mempertemukan telaah historis-filologis klasik dengan agenda komputasional-digital, lewat perangkuman periodisasi, transmisi, dan pemalsuan ke dalam satu narasi yang padu.\footnote{Tanjung, "The Role of Tadwin Hadith."; Abdulrahman, "The Future of Hadith Studies in the Digital Age."; Mosa, "Synergizing Structure and Semantics.".} Berbeda dari riset satu-fase --- misalnya artikulasi wibawa kenabian pada \textit{al-Muwatta'} serta \textit{Musannaf}\footnote{Musadad, "Articulating the Prophetic Authority.".}, atau polemik kesakralan Madinah\footnote{Kara, "Debating the Origins.".} --- riset ini memetakan lima kurun: Nabi, sahabat, tabi'in, ulama, dan modern, sebagai satu bingkai analitis bagi perkembangan metodologi transmisi beserta kritik.\footnote{Razak, "Evolution of Hadith Textual Criticism."; Lutfianto, "History of Hadith Writing.".} Lain halnya dengan riset komputasional yang teknosentris\footnote{Alghamdi, "Pretrained Models Against Traditional Machine."; Shaaban, "Prophetic Authorial Style Modeling.".}, riset ini mendudukkan teknologi selaku penerus etos \textit{tathabbut}, \textit{rihlah}, serta \textit{jarh-ta'dil}, bukan substitusi otoritas ulama.\footnote{Jaiyeoba, "Hadith Preservation."; Muhamed Ali, "Al-jarh Wa Al-ta'dil (criticism and Praise)."; Gil, "Tracing the Rihla in Buldan Books.".} Sumbangsihnya mencakup pemadatan historiografis peralihan lisan-tulisan-kodifikasi-digital, taksonomi pemicu pemalsuan berikut bangunan penyelamatannya hingga zaman algoritmik, serta agenda verifikasi digital berlandas paradigma \textit{maqasid}.\footnote{Bin Jamil, "Between Traditionalising and Futuring."; Abdulrahman, "The Future of Hadith Studies in the Digital Age."; Kamran, "Semantichadith.".} Tulisan ini diharapkan menjadi gerbang yang layak menuju pembahasan periodisasi, penghimpunan, dan pemalsuan hadis pada Bab II, sekaligus mengukuhkan relevansi ilmu hadis selaku disiplin hidup di tengah pergantian zaman.\footnote{Tanjung, "The Role of Tadwin Hadith."; Razak, "Evolution of Hadith Textual Criticism."; Abdullah, "Viralitas Dan Validitas.".}

\phantomsection
\subsection*{B. Rumusan Masalah}
\addcontentsline{toc}{subsection}{B. Rumusan Masalah}

Berdasarkan latar belakang di atas, rumusan masalah penelitian ini adalah:

\begin{enumerate}[leftmargin=1.27cm,itemsep=2pt,parsep=2pt]
\item Bagaimana periodisasi sejarah perkembangan hadis dari masa Nabi hingga era modern?

\item Bagaimana proses penghafalan dan penghimpunan hadis dari tradisi lisan hingga kodifikasi resmi?

\item Apa faktor-faktor yang mendorong pemalsuan hadis dan bagaimana upaya penyelamatannya?

\end{enumerate}
\phantomsection
\subsection*{C. Tujuan Penulisan}
\addcontentsline{toc}{subsection}{C. Tujuan Penulisan}

Sejalan dengan rumusan masalah, penelitian ini bertujuan:

\begin{enumerate}[leftmargin=1.27cm,itemsep=2pt,parsep=2pt]
\item Mengetahui dan menganalisis periodisasi sejarah perkembangan hadis.

\item Memahami proses penghafalan dan penghimpunan hadis.

\item Menjelaskan faktor-faktor pemalsuan hadis serta upaya penyelamatan yang dilakukan ulama.

\end{enumerate}
\phantomsection
\subsection*{D. Manfaat Penulisan}
\addcontentsline{toc}{subsection}{D. Manfaat Penulisan}

\begin{itemize}[leftmargin=1.27cm,itemsep=2pt,parsep=2pt]
\item Manfaat teoretis: menambah kekayaan khazanah ilmu hadis melalui perangkuman historiografis atas temuan klasik dan mutakhir, utamanya transisi epistemologis dari kelisanan menuju pembukuan serta digitalisasi.\footnote{Tanjung, "The Role of Tadwin Hadith."; Razak, "Evolution of Hadith Textual Criticism."; Abdulrahman, "The Future of Hadith Studies in the Digital Age.".}

\item Manfaat praktis: menyediakan bahan literasi bagi mahasiswa, pengajar, dan publik untuk memeriksa hadis di lingkungan digital, agar terhindar dari \textit{maudu'} sekaligus mampu memakai perangkat verifikasi \textit{takhrij} digital secara kritis.\footnote{Abdullah, "Viralitas Dan Validitas."; Supriyadi, "Action Research in Hadith Literacy."; Syukur, "Living Hadith in the Age of Artificial Intelligence.".}

\end{itemize}
\phantomsection
\subsection*{E. Metode Penulisan}
\addcontentsline{toc}{subsection}{E. Metode Penulisan}

Penelitian ini memakai telaah kepustakaan (\textit{library research}) yang bersifat kualitatif-deskriptif dan historis-analitis. Data meliputi 35 rujukan utama dari jurnal bereputasi: riset periodisasi\footnote{Musadad, "Articulating the Prophetic Authority."; Tanjung, "The Role of Tadwin Hadith."; Razak, "Evolution of Hadith Textual Criticism.".}, penghimpunan beserta preservasi\footnote{Rashwan, "Hadith As Oral Literature Through."; Lutfianto, "History of Hadith Writing."; Jaiyeoba, "Hadith Preservation."; Obeid, "Al-bukhari's Sources.".}, kritik berikut pemalsuan\footnote{Muhamed Ali, "Al-jarh Wa Al-ta'dil (criticism and Praise)."; Su, "'ali B. Al-madini (161-234/778-849)."; Abbas, "The Concept of Munkar."; Mohamed, "Linguistic Features Evaluation for Hadith.".}, dan kajian digital-kontemporer.\footnote{Abdulrahman, "The Future of Hadith Studies in the Digital Age."; Matin Bin Salman, "Reconstructing Hadith Discourse in the Digital Age."; Alghamdi, "Pretrained Models Against Traditional Machine."; Abdullah, "Viralitas Dan Validitas.".} Data dihimpun lewat kajian dokumentasi dan perangkuman literatur; analisisnya memakai analisis isi melalui reduksi, pengelompokan tematik merujuk rumusan masalah, dan penarikan simpulan deduktif-analitis.

